\documentclass[10pt,a4paper]{amsart}
\usepackage{amsmath,amssymb,amsthm,graphicx}
% Retained preamble for the Stacks Project chapter {CH} (shared project preamble).
% The chapter file pinned-original-source.tex only carries \input{{preamble}} on its first line,
% so no preamble text lives inside the chapter itself.  The Stacks Project uses one single shared
% preamble for every chapter: the snapshot stages it three times (sources/stacks-chow/preamble.tex,
% sources/stacks-etale-cohomology/preamble.tex and focused-seeds/AGP/sources/stacks-snapshot-20261011/
% preamble.tex) and those copies are identical in size.  The block below is copied mechanically,
% line for line, from sources/stacks-chow/preamble.tex lines 14-18 and 154-217: the Xy-pic setup and
% the native theorem declarations together with the mathematical macros.  Nothing was added, removed
% or redefined; the wrapper class and the amsmath/amssymb/amsthm/graphicx base packages are supplied
% by the assembler.
% For commutative diagrams we use Xy-pic
% Chapter: stacks-spaces-more-morphisms.tex -- More on Morphisms of Spaces.
\usepackage[all]{xy}

% We use 2cell for 2-commutative diagrams.
\xyoption{2cell}

% Theorem environments.
%
\theoremstyle{plain}
\newtheorem{theorem}[subsection]{Theorem}
\newtheorem{proposition}[subsection]{Proposition}
\newtheorem{lemma}[subsection]{Lemma}

\theoremstyle{definition}
\newtheorem{definition}[subsection]{Definition}
\newtheorem{example}[subsection]{Example}
\newtheorem{exercise}[subsection]{Exercise}
\newtheorem{situation}[subsection]{Situation}

\theoremstyle{remark}
\newtheorem{remark}[subsection]{Remark}
\newtheorem{remarks}[subsection]{Remarks}

\numberwithin{equation}{subsection}

% Macros
%
\def\lim{\mathop{\mathrm{lim}}\nolimits}
\def\colim{\mathop{\mathrm{colim}}\nolimits}
\def\Spec{\mathop{\mathrm{Spec}}}
\def\Hom{\mathop{\mathrm{Hom}}\nolimits}
\def\Ext{\mathop{\mathrm{Ext}}\nolimits}
\def\SheafHom{\mathop{\mathcal{H}\!\mathit{om}}\nolimits}
\def\SheafExt{\mathop{\mathcal{E}\!\mathit{xt}}\nolimits}
\def\Sch{\mathit{Sch}}
\def\Mor{\mathop{\mathrm{Mor}}\nolimits}
\def\Ob{\mathop{\mathrm{Ob}}\nolimits}
\def\Sh{\mathop{\mathit{Sh}}\nolimits}
\def\NL{\mathop{N\!L}\nolimits}
\def\CH{\mathop{\mathrm{CH}}\nolimits}
\def\proetale{{pro\text{-}\acute{e}tale}}
\def\etale{{\acute{e}tale}}
\def\QCoh{\mathit{QCoh}}
\def\Ker{\mathop{\mathrm{Ker}}}
\def\Im{\mathop{\mathrm{Im}}}
\def\Coker{\mathop{\mathrm{Coker}}}
\def\Coim{\mathop{\mathrm{Coim}}}

% Boxtimes
%
\DeclareMathSymbol{\boxtimes}{\mathbin}{AMSa}{"02}

%
% Macros for moduli stacks/spaces
%
\def\QCohstack{\mathcal{QC}\!\mathit{oh}}
\def\Cohstack{\mathcal{C}\!\mathit{oh}}
\def\Spacesstack{\mathcal{S}\!\mathit{paces}}
\def\Quotfunctor{\mathrm{Quot}}
\def\Hilbfunctor{\mathrm{Hilb}}
\def\Curvesstack{\mathcal{C}\!\mathit{urves}}
\def\Polarizedstack{\mathcal{P}\!\mathit{olarized}}
\def\Complexesstack{\mathcal{C}\!\mathit{omplexes}}
% \Pic is the operator that assigns to X its picard group, usage \Pic(X)
% \Picardstack_{X/B} denotes the Picard stack of X over B
% \Picardfunctor_{X/B} denotes the Picard functor of X over B
\def\Pic{\mathop{\mathrm{Pic}}\nolimits}
\def\Picardstack{\mathcal{P}\!\mathit{ic}}
\def\Picardfunctor{\mathrm{Pic}}
\def\Deformationcategory{\mathcal{D}\!\mathit{ef}}
\newenvironment{reference}{}{}
\newenvironment{slogan}{}{}
\newenvironment{history}{}{}
\usepackage{cleveref}
\crefname{originallemma}{original lemma}{original lemmas}
\Crefname{originallemma}{Original lemma}{Original lemmas}
\makeatletter
\newlabel{definition-lci}{{069F}{source}{}{sourcepointer-1}{}}
\newlabel{definition-lci@cref}{{[originallemma][0][]069F}{[1][0][]source}{}{}{}}
\makeatother
\crefname{originallemma}{original lemma}{original lemmas}
\Crefname{originallemma}{Original lemma}{Original lemmas}
\makeatletter
\newlabel{more-morphisms-lemma-base-change-lci-fibres}{{06B8}{source}{}{sourcepointer-2}{}}
\newlabel{more-morphisms-lemma-base-change-lci-fibres@cref}{{[originallemma][0][]06B8}{[1][0][]source}{}{}{}}
\makeatother
\crefname{originallemma}{original lemma}{original lemmas}
\Crefname{originallemma}{Original lemma}{Original lemmas}
\makeatletter
\newlabel{spaces-properties-lemma-etale-open}{{03IR}{source}{}{sourcepointer-3}{}}
\newlabel{spaces-properties-lemma-etale-open@cref}{{[originallemma][0][]03IR}{[1][0][]source}{}{}{}}
\makeatother
\crefname{originallemma}{original lemma}{original lemmas}
\Crefname{originallemma}{Original lemma}{Original lemmas}
\makeatletter
\newlabel{spaces-properties-lemma-open-subspaces}{{03BZ}{source}{}{sourcepointer-4}{}}
\newlabel{spaces-properties-lemma-open-subspaces@cref}{{[originallemma][0][]03BZ}{[1][0][]source}{}{}{}}
\makeatother
\title{The Koszul locus of a map over a flat finitely presented base family is open and stable under base change (spaces)}
\author{The Stacks Project Authors}
\date{stacks-project:spaces-more-morphisms.tex:local-snapshot-20261011; GNU FDL 1.2 or later}
\begin{document}
\maketitle
\noindent\textit{Editorial note.} Statement span is lines 11382-11398 of sources/stacks-spaces-more-morphisms/pinned-original-source.tex; proof span is lines 11400-11438, the complete author proof (37 body lines).\par\medskip
\noindent\textit{Editorial note.} Navigation stubs: 4 cited-but-not-included labels are supplied as non-load-bearing pointers with their pinned tags where the snapshot tag index has them.\par\medskip
\noindent\textit{Editorial note.} extra\_packages is empty on purpose: the retained preamble already loads xy with options and a second declaration clashes in the assembler wrapper.\par\medskip
\noindent\textit{Editorial note.} Text is retained verbatim, including the author's own typos and placeholders.\par\medskip
\noindent\textit{Editorial note.} Counted B side (both complete); schemes twin is imported as the load-bearing input, so its own label is the uncounted side of this pair per ruling v3.\par\medskip
\noindent\textit{Allowed standard prerequisites (editorial).}\begin{itemize}
\item Explicit prerequisite boundary. More on Morphisms, lemma base-change-lci-fibres (schemes twin); chapter definition-lci; Properties of Spaces: open-subspaces, etale-open.
\end{itemize}\medskip
\ifdefined\hypertarget\hypertarget{sourcepointer-1}{}\fi
\noindent\textit{Source reference.} Original lemma 069F: sources/stacks-spaces-more-morphisms/pinned-original-source.tex, source lines 11153--11164. Cited by the retained statement or proof but not part of the collected spans; supplied only as a navigation pointer. The argument does not depend on its content.\par\medskip
\ifdefined\hypertarget\hypertarget{sourcepointer-2}{}\fi
\noindent\textit{Source reference.} Original lemma 06B8: sources/stacks-spaces-more-morphisms/pinned-original-source.tex, source lines 11438--11438. Cited by the retained statement or proof but not part of the collected spans; supplied only as a navigation pointer. The argument does not depend on its content.\par\medskip
\ifdefined\hypertarget\hypertarget{sourcepointer-3}{}\fi
\noindent\textit{Source reference.} Original lemma 03IR: sources/stacks-spaces-more-morphisms/pinned-original-source.tex, source lines 11438--11438. Cited by the retained statement or proof but not part of the collected spans; supplied only as a navigation pointer. The argument does not depend on its content.\par\medskip
\ifdefined\hypertarget\hypertarget{sourcepointer-4}{}\fi
\noindent\textit{Source reference.} Original lemma 03BZ: sources/stacks-spaces-more-morphisms/pinned-original-source.tex, source lines 11438--11438. Cited by the retained statement or proof but not part of the collected spans; supplied only as a navigation pointer. The argument does not depend on its content.\par\medskip
\begin{lemma}
\label{lemma-base-change-lci-fibres}
Let $S$ be a scheme. Consider a commutative diagram
$$
\xymatrix{
X \ar[rr]_f \ar[rd]_p & & Y \ar[ld]^q \\
& Z
}
$$
of algebraic spaces over $S$. Assume that both $p$ and $q$
are flat and locally of finite presentation.
Then there exists an open subspace $U(f) \subset X$
such that $|U(f)| \subset |X|$ is the set of points where $f$ is Koszul.
Moreover, for any morphism of algebraic spaces $Z' \to Z$, if
$f' : X' \to Y'$ is the base change of $f$ by $Z' \to Z$, then
$U(f')$ is the inverse image of $U(f)$ under the projection $X' \to X$.
\end{lemma}

\begin{proof}
This lemma is the analogue of
More on Morphisms, Lemma \ref{more-morphisms-lemma-base-change-lci-fibres}
and in fact we will deduce the lemma from it. By
Definition \ref{definition-lci}
the set $\{x \in |X| : f \text{ is Koszul at }x\}$ is
open in $|X|$ hence by
Properties of Spaces, Lemma \ref{spaces-properties-lemma-open-subspaces}
it corresponds to an open subspace $U(f)$ of $X$. Hence we only need to
prove the final statement.

\medskip\noindent
Choose a scheme $W$ and a surjective \'etale morphism $W \to Z$.
Choose a scheme $V$ and a surjective \'etale morphism $V \to W \times_Z Y$.
Choose a scheme $U$ and a surjective \'etale morphism $U \to V \times_Y X$.
Finally, choose a scheme $W'$ and a surjective \'etale morphism
$W' \to W \times_Z Z'$.
Set $V' = W' \times_W V$ and $U' = W' \times_W U$, so that we obtain
surjective \'etale morphisms $V' \to Y'$ and $U' \to X'$.
We will use without further mention an \'etale morphism of algebraic spaces
induces an open map of associated topological spaces (see
Properties of Spaces, Lemma
\ref{spaces-properties-lemma-etale-open}).
Note that by definition $U(f)$ is the image in $|X|$ of the set $T$
of points in $U$ where the morphism of schemes $U \to V$ is Koszul.
Similarly, $U(f')$ is the image in $|X'|$ of the set $T'$ of points in
$U'$ where the morphism of schemes $U' \to V'$ is Koszul. Now, by construction
the diagram
$$
\xymatrix{
U' \ar[r] \ar[d] & U \ar[d] \\
V' \ar[r] & V
}
$$
is cartesian (in the category of schemes). Hence the aforementioned
More on Morphisms, Lemma \ref{more-morphisms-lemma-base-change-lci-fibres}
applies to show that $T'$ is the inverse image of $T$. Since
$|U'| \to |X'|$ is surjective this implies the lemma.
\end{proof}

\end{document}
